\documentclass[a4paper,11pt]{article}

\usepackage{revy}
\usepackage{amssymb}
\usepackage{amsmath}
\usepackage{commath}
\usepackage{amsfonts}

\revyname{MatematikRevyen}
\revyyear{2026}
\version{1}
\eta{2-3 minutter}
\status{Færdig}

\title{ChatPhD}
\author{Malthe '20}

\begin{document}
\maketitle

\begin{roles}
\role{S1} Studerende 
\role{S2} Studerende
\role{P} PhD-studerende


\end{roles}

\begin{props}
    \prop{2} Stole    
    \prop{1} Bord
    \prop{2} Laptops
\end{props}


\begin{sketch}

\scene To studerende sidder ved siden af hinanden, begge med deres laptops fremme. S1 kigger over på S2's skærm og ryster på hovedet. 

\says{S1} Bro, bruger du stadig ChatGPT? Det er der altså ingen der bruger længere. 

\says{S2} Nå, er der kommet en ny version af Claude eller hvad?

\says{S1} Nej nej nej, jeg bruger sgu da ChatPhD.

\says{S2} ChatPhD? Det har jeg altså aldrig hørt om. 

\says{S1} Det er KUs egen interne model. Den kan svare på alle matematiske spørgsmål, og den kender til alle dine kurser. Det er lige som at have din egen personlige instruktor.

\says{S2} Wow, jeg vidste slet ikke at vi havde vores egen model her. Men det må have været dyrt at skabe sådan én. Skal jeg så betale for at bruge den?

\says{S1} Det har selvfølgelig kostet både tid og penge at træne den, men her er det smarte trick: Man har trænet den på instituttets egne kurser, som jo alligevel ville blive afholdt, så selvom det har taget nogle år at træne den, så har det været ret billigt. Som KU-studerende har du derfor gratis adgang til at bruge den. Du skal bare banke på en tilfældig dør i E-bygningen for at prompte den. 

\says{S2} Det overrasker mig faktisk lidt at KU har lavet deres egen model. Der er jo alle klimamålene, og datacentre fylder vildt meget og sluger både strøm og vand.

\says{S1} Ja, den fylder godt nok noget. Man har været nødt til at inddrage hele kælderetagen på HCØ til at hoste den, men så ved du til gengæld også at alting er lagret lokalt. Men til gengæld har jeg hørt, at den har et minimalt strømforbrug, hvoraf størstedelen kommer når lyset ikke bliver slukket henover weekenden. Det er så meget bedre end andre AIs. \\
Vandforbruget er nok mere sammenligneligt med andre modeller, men det er hovedsageligt på grund af et absurd højt forbrug af kaffe.

\says{S1} Okay, fedt nok, jeg skal helt klart prøve den. Kan den lave mine afleveringer for mig?

\says{S2} Endnu bedre. Den har et indbygget etisk kodeks, som sikrer sig at du ikke begår eksamenssnyd. Den styrer dig udenom om alle fælderne for plagiat, og den guider dig i retningen af den rigtige løsning, samtidig med at den sørger for, at du lærer mest muligt. 

\says{S1} Det lyder voldsomt avanceret. Hvordan er det lykkedes KU at lave en model der er så meget bedre end tech-giganternes?

\says{S2} De udnytter et splinternyt agent-system, hvor de har mange forskellige agenter med forskellige specialiseringer, der både kan arbejde uafhængigt af hinanden og snakke sammen. Og det bedste er, at når der er ikke er nogen, der aktivt prompter den, så kører den faktisk stadig i baggrunden hele døgnet, alle årets dage, også hen over juleferien, og alle agenterne bliver bare klogere og klogere hele tiden.

\says{S2} Det lyder altså næsten lidt skræmmende. Er du ikke bange for at den bliver for klog og prøver at overtage verden?

\says{S1} Jeg har hørt rygter om enkelte agenter, der har været tæt på at stikke af, men jeg tror altså ikke for alvor at ChatPhD kommer til at udvikle bevidsthed.

\scene P kommer forvirret ind på scenen med uglet hår og dybe render under øjnene. Hun prøver at dække for det skarpe lys med hænderne.

\says{P} Hvor er mit kontor? Hvilken dag er det?

\scene S1 og S2 rejser sig og fører forsigtigt P væk fra scenen. 

\says{S2} Kom min ven, tilbage i kælderen med dig.  

\says{S1} Det er altså derfor man skal huske at lukke vinduerne når man går fra HCØ.



\scene Lys ned






\end{sketch}

\end{document}

